% Combined from chapters/quotes/*.tex
% Merge command: python3 convert_peing_qa_to_latex.py --combine

\chapter*{問答語錄}
\addcontentsline{toc}{chapter}{問答語錄}

\chapter[今まで誰にも言ってない秘密を教えて！]{今まで誰にも言ってない秘密を教えて！}
\qaanswer{
言わないよ。秘密は言わないこそ秘密だよ。
}

\chapter[您的吳語播音體建構是否已有大概時間表，例如何時開始發表中篇及長篇的播音體文章？]{您的吳語播音體建構是否已有大概時間表，例如何時開始發表中篇及長篇的播音體文章？}
\qaanswer{
吳語播音體不是我建構的，我只是參與建構的人和目前的對外發信人。播音體建構大約始於2009年，10年下來語法體系和構詞法已經建構的差不多了。中篇、長篇，各種文體也都嘗試過。現在主要的精力是編寫『簡明漢吳辭典』、最佳化吳小字的字型以及定期發布『吳語播音體語法和寫作指導』，讓更多的人能夠接觸和學習吳語播音體。
}

\chapter[您所整理的吳語播音體語法字尾主要適用的地區是？將來能覆蓋到整個吳語區嗎？]{您所整理的吳語播音體語法字尾主要適用的地區是？將來能覆蓋到整個吳語區嗎？}
\qaanswer{
基本上適用於吳語、徽語多數地區，是涵蓋南北吳大框架下的語法建構。透過建構，重新梳理了漢字圈各語言的語法脈絡，通宵吳語播音體語法的人甚至可以在此基礎上幫助老湘語、淮東語、贛語、膠遼語甚至漢語普通話進行高階語法建構。題外話，漢語普通話需要嚴謹嚴密的語法梳理，否則這門語言對法學、神學、哲學及自然科學研究是有障礙的。
}

\chapter[你老能否闡釋一下環東海文明圈的概念，大概包括哪些地區？]{你老能否闡釋一下環東海文明圈的概念，大概包括哪些地區？}
\qaanswer{
雖然我不喜歡用「自古以來」這個詞，但是受到「東中國沿岸寒流」「黑潮」、「臺灣暖流」、「對馬暖流」的季節性影響，東海自古以來就是一條有時限的海上高速公路。它連線了：山東半島、遼東半島、日本列島、朝鮮半島(高句麗、新羅、百濟)、沖繩(琉球)、臺灣(夷洲)、菲律賓(呂宋)。

這些地區的貿易據點構成的網路至少在4世紀初就成型了。貿易中涉及的物資、人員交往帶來了文化、文明的相互影響。尤其是受到東亞大陸政權影響相對更弱的日本、新羅、百濟、夷洲、呂宋與吳越的關係更為緊密。

從4世紀後半葉到10世紀末自有一套政治、經濟體系相對獨立的運作著並作用到以上各方。11世紀到14世紀由於宋帝國、元帝國先後對遠洋航海貿易的支援，這種關係得到了加強，先有南宋時期的宋日貿易，後有元帝國的「世界貿易」都是環東海文明圈的巔峰時期。之後由盛轉衰，進入到明帝國前期鄭和七下西洋後是一個長達近600年的漫長的海禁期，環東海文明戛然而止，少量文明交流的碎片被吳越武裝走私商、明政府軍、西葡荷商團僱傭軍、日本南部諸藩相互間的角力所裹挾，正常的大範圍的交流被明帝國的政策幹擾，直至清帝國時期全面中止。
}

\chapter[您怎麼看待杭州話的一些官話特徵？]{您怎麼看待杭州話的一些官話特徵？}
\qaanswer{
這個問題太空泛了，很難回答。是問我杭州話的形成脈絡呢？還是問我對杭州話帶有官話特徵的態度？如果是後者我只想說，其實吳越老府城的口語都官，事實上蘇白除了吳語腔調本身就是官話語法，沒比杭州話的官話特徵要少到哪裡去。
}

\chapter[江淮官話是官話化的吳語嗎？]{江淮官話是官話化的吳語嗎？}
\qaanswer{
這個我不敢瞎欽點。江淮官話內部差異大了海了。得分片討論，和吳語最接近的是淮東語通泰片，共享海量詞彙集和部分北吳語法標記。
}

\chapter[您知道有誰系統地寫了吳越國史嗎？如果沒有，我打算在墻內現擬定一份初稿。因為才疏學淺，期間遇到問題，會繼續在此處向您發問。]{您知道有誰系統地寫了吳越國史嗎？如果沒有，我打算在墻內現擬定一份初稿。因為才疏學淺，期間遇到問題，會繼續在此處向您發問。}
\qaanswer{
吳語公眾號有人在更吧。如果在牆內建議敏感內容還是不要公佈的好。

寫史先得確立歷史觀，構建吳越史採用的歷史觀建議參考UK的歷史構建觀。因為吳越和UK非常類似，同樣都是由少數原住民，歷史上不斷地以武力、半武力方式遷入的移民層層交錯疊合而成的文明體；同樣在在出現穩定的文化、文明特質後（吳越是南宋以後\ruby{󱚾󱤦}{UK}是諾曼征服者進入以後），都沒有經歷過釜底抽薪的大浩劫（戰爭、瘟疫引起的人口滅絕、人口置換），且內部比較好的融合出一套獨具自身特點的不同與周邊地區的語言、文化、風俗及工商產業結構。
}

\chapter[上海市區話已經無法搶救，是否能從 松江 嘉定，浦東崇明等地有殘留原生詞彙語法的重新構建出一套出來？這些郊縣，那個語法標記詞彙儲存最完整？]{上海市區話已經無法搶救，是否能從 松江 嘉定，浦東崇明等地有殘留原生詞彙語法的重新構建出一套出來？這些郊縣，那個語法標記詞彙儲存最完整？}
\qaanswer{
上海市區話本身在形成過程中就一種藍青官話。因為本身就沒有標準，語法體系又都是官話一脈的基本上沒辦法拯救，放到五十年前也救不了。上海郊縣語法標記保留的都還可以，特別是松江，有很多熱心的同胞在整理和紀錄語法、詞彙，保留了大量的文獻記載。
}

\chapter[我知道吳語小字和你們構建的吳語書面語是兩回事，書面語和語言的思維，是遠遠比表音字型重要得多。 但是，在你們團隊建構的吳語字典和寫作指南尚未出爐之前，大部分的人都是吳語初學者吧？更何況，現在連吳語的發音腔調都正在被官話置換當中。 問題來了：那些被你認為官話成分極高的吳語方言（城裡的），這些吳語方言的作品（歌曲或影片） ，值得被標上吳語小字的字幕嗎？ ]{我知道吳語小字和你們構建的吳語書面語是兩回事，書面語和語言的思維，是遠遠比表音字型重要得多。 \\[0.35em] 但是，在你們團隊建構的吳語字典和寫作指南尚未出爐之前，大部分的人都是吳語初學者吧？更何況，現在連吳語的發音腔調都正在被官話置換當中。 \\[0.35em] 問題來了：那些被你認為官話成分極高的吳語方言，也就是城裡頭的，這些吳語方言的作品，譬如歌曲或影片 ，值得被標上吳語小字的字幕嗎？ \\}
\qaanswer{
使用吳小字純屬個人自由，只要妳開心可以用來標註任何語言的發音。研發\ruby{}{iphone}硬體的人總不能規定使用者說「必須這麼用，不能那麼用」吧？
}

\chapter[請問現行從思源宋體改的小字，製作者從零到完成 總共花多少時間? 另外，找專業的字型公司造字，大概需要花多少錢？(日幣)...... 加油，祝早日完成!! \textasciitilde{}\textasciitilde{}]{請問現行從思源宋體改的小字，製作者從零到完成 總共花多少時間? \\[0.35em] 另外，找專業的字型公司造字，大概需要花多少錢(日幣)...... \\[0.35em] 加油，祝早日完成!! \textasciitilde{}\textasciitilde{}}
\qaanswer{
思源宋體改的吳小字從零開始到第一版出來字型+字元加在一起花了7個月時間。之後有了經驗就快了，最新版的只花了1個月時間。第一版和最新版差異非常大。 問日本的字型公司造字需要先做好字元項的準備，需要具體詢價。目前在試做黑體。
}

\chapter[請問...目前你們在自製的黑體的過程中， 具體會用到哪些軟體?推薦用哪些軟體 對新手比較好操作? 你知道的]{請問...目前你們在自製的黑體的過程中， \\[0.35em] 具體會用到哪些軟體?推薦用哪些軟體 對新手比較好操作? \\[0.35em] 你知道的}
\qaanswer{
做過字元和字型的朋友知道，雖然有專業軟體(迫真)輔助，字元字型要做的好看那是一門玄學。事實上目前就我一個人抱著玩耍的心態在做吳小字的更紗黑體。軟體用的是FontCreator。
}

\chapter[也不算問題吧。我只在網上學過國際音標和一點基礎語音學，而您的和其他人的吳語拼音教程沒有標註寬式國際音標，而我本地是最西部的太湖片地區與太湖片核心地區很多具體字音不一樣，所以遇到一些學習困難。另外，我覺得您應該推廣（或自行設計）更貼近英語拼讀規律的吳語拼音，比如「zh.wikipedia.org/wiki/上海話拉丁化方案」上的「江拉、滬一、錢一」。]{也不算問題吧。我只在網上學過國際音標和一點基礎語音學，而您的和其他人的吳語拼音教程沒有標註寬式國際音標，而我本地是最西部的太湖片地區與太湖片核心地區很多具體字音不一樣，所以遇到一些學習困難。另外，我覺得您應該推廣（或自行設計）更貼近英語拼讀規律的吳語拼音，比如「zh.wikipedia.org/wiki/上海話拉丁化方案」上的「江拉、滬一、錢一」。}
\qaanswer{
妳提出的是一個很多沒接觸過吳語拼音的朋友會問的典型問題

首先吳語拼音不是音值拼音，而是音位拼音，所謂音位指的是一個音位對應多個音值，比如同一個音位韻母eu，在吳語區各地的音值是不同的，這樣方便各地人按自己的習慣讀音，不需要人為的規定什麼標準音(定標準音又是一個誰也不服誰的口水仗)。

如果按國際音標來，那麼實際情況下吳語區一個鎮一個村就是一個拼法，整套系統下來估計會是天文數字吧！而且國際音標學習成本極高。事實上英語就是音位拼音語言，同一個拼寫英語區各地實際讀音都不太一樣。

吳語拼音是目前我個人覺得最科學也是最便於學習的一套拼音方案。希望我的回答能讓妳滿意。
}

\chapter[請問，歷史上曾經有從上古或中古漢語分化出的齊語、楚語和蜀語嗎？]{請問，歷史上曾經有從上古或中古漢語分化出的齊語、楚語和蜀語嗎？}
\qaanswer{
這個我不清楚，不過我知道以下3個關於四川地區語言變化脈絡的知識點：

1.在北宋中葉有士大夫的文獻中記載，蜀語(梁益方言)類閩。也就是說當時的蜀語和當時閩地的語言有很多近似的地方。但北宋時的蜀語和現在的四川話差異很大，整個面貌完全不同。這涉及到宋代到清初四川地區的數次人口置換。

2.現在的四川話顯然是清中期人口佔絕對優勢的兩湖人、贛人填四川時覆蓋過去的西南官話和少數碎片化了的老湘語、贛語。它們與之前傳入四川的江淮軍話，殘餘的老蜀語一起混合穩定後顯了現在西南官話四川區片的狀態。

3.在清以前的明代整個西南地區的各殖民點大城就已經被作為明初軍事征服者的朱元璋的江淮部隊「軍話」(江淮官話)和徽州移民講的徽語(高地吳語的一種)覆蓋過一次了。

2、3的情況導致大的西南官話區的形成，現在四川、貴州、雲南、湖北、湖南的語言互通率很高。而且越是晚期被漢地風俗所同化的少數民族地區西南官話使用率反而越高。比如湘黔交界的定居苗族不使用老湘語反而使用西南官話。雲南、貴州很多土司地區的情況也類似，越晚被同化的地區講的官話反而越「地道」。
}

\chapter[請問，您怎麼看待日本的漢字簡化？]{請問，您怎麼看待日本的漢字簡化？}
\qaanswer{
我個人的意見是沒有意見。

人家按照自己的經驗主義進行簡化那是人家的事。這符合人家實際的需要。
}

\chapter[吳語公眾號上廣告一本今年出版的書 書名叫做 吳語方言學 這本值得買嗎]{吳語公眾號上廣告一本今年出版的書 \\[0.35em] 書名叫做 吳語方言學 \\[0.35em] 這本值得買嗎}
\qaanswer{
靠一本書就能學懂不太現實。也別相信書本上的每個字。但是如果想要瞭解吳語，需要多看各類關於吳語的論文。吳語研究現在的問題是一群不會吳語的官話學者掌握了吳語研究的話語權，而另一些以吳語為母語的研究者的立場並沒有把吳語當作一門獨立於漢語普通話的語言。這個立場決定了他們都不會為構建吳語的書面化而出力，事實上他們做的吳語研究只能使這門語言逐步形骸化，化石化。吳語如果不能被寫出來，那麼這門語言在很短的時間內就會消亡。
}

\chapter[有人覺得蘇錫常和皖南吳語區的散沙化和對母語的不熟悉比浙江嚴重很多，尤其是在青少年一代中，您贊同這個觀點嗎]{有人覺得蘇錫常和皖南吳語區的散沙化和對母語的不熟悉比浙江嚴重很多，尤其是在青少年一代中，您贊同這個觀點嗎}
\qaanswer{
我不贊同這樣的說法。因為這些都無法量化並被統計出來。無論是蘇南、徽州還是浙江都面臨普通話的強力影響與世代置換。而且在我們這代人可見的未來，吳語徹底滅亡。討論誰比誰對母語更不熟悉跟討論用幾種方法殺死將要被烹飪的章魚一樣，對與受體而言意義不大。我倒是對「如何讓吳語延續下去」，「章魚如何逃脫廚師的魔爪回到大海」這類的議題更感興趣。
}

\chapter[母語詞彙可以從別的語言借用，這叫做外來語， 這是很常見的現象。 那構詞法和文法也可以從別的語言借過來，並紮根於母語中嗎？]{母語詞彙可以從別的語言借用，這叫做外來語， \\[0.35em] 這是很常見的現象。 \\[0.35em] 那構詞法和文法也可以從別的語言借過來，並紮根於母語中嗎？}
\qaanswer{
最穩妥的方式是按照經驗主義的方式來構建，也就是說外來語的文法、構詞法首先要能被吳語區的人在使用時普遍接受，但同時又不能與吳語本有的文法、構詞法產生邏輯矛盾。比如吳語非常講究區別詞性，而漢語就很難區分詞性了。基本上是中古漢語、近古漢語、梵語、中古波斯語、契丹語的一些詞彙作為詞根融入到了吳語中。直接從外來語拿來構詞法和文法我相信很難有人認可，也不知如何來使用他。現在主要的問題是吳語區的人掌握的吳語基本也都是些被官話、普通話嚴重滲透失去吳語語法特點的藍青官話。按我的標準這些都算不上吳語，是一種帶著吳語口音、偶爾有吳語詞彙的「普通話」。這些人對於吳語語法的瞭解、構詞法的瞭解基本上是一竅不通的。建構的前提是瞭解自身。與其一開始考慮借鑑別人不如先好好了解自己。懂得自身才能知道外界哪些是需要的哪些是不需要的。
}

\chapter[吳越人在日韓有組織度比較高的同鄉會這樣的組織嗎?]{吳越人在日韓有組織度比較高的同鄉會這樣的組織嗎?}
\qaanswer{
就現在這個時間點來說基本上不存在。
}

\chapter[您老有計劃做原生吳語的比較系統的發音、文法教學影片嗎？]{您老有計劃做原生吳語的比較系統的發音、文法教學影片嗎？}
\qaanswer{
文法一直有出啊，另外吳語沒有所謂標準音，所以吳語拼音這種一個音位對應多個音值的拼音方案是最佳選。人為定標準音那是大一統帝國思維，都2019年了，觀念是不是得轉變一下？經驗主義的思路是使用某種語言的有教養的人群會自然形成一種語音腔調方式，這種腔調就會成為該語言的標準音。吳語的高階化路還很漫長。
}

\chapter[普通話能整理出文法嗎？普通話當然有文法呀！但是在這年頭，只能在正經的文章才能看見正常的文字了。網路上的用詞構句很散漫，新聞媒體則是精通語言，卻故意玩弄文字，而我竟習慣這些混亂的用法而不自知。怪我自己不愛讀書用腦哩⋯⋯不只你這樣批評過中文（普通話？）的現象，我聽過至少二十個立場相左、背景不同的人，都對此抱著相同的看法：現在的中文也好、普通話也罷，是很混亂的。]{普通話能整理出文法嗎？普通話當然有文法呀！但是在這年頭，只能在正經的文章才能看見正常的文字了。網路上的用詞構句很散漫，新聞媒體則是精通語言，卻故意玩弄文字，而我竟習慣這些混亂的用法而不自知。怪我自己不愛讀書用腦哩⋯⋯不只你這樣批評過中文（普通話？）的現象，我聽過至少二十個立場相左、背景不同的人，都對此抱著相同的看法：現在的中文也好、普通話也罷，是很混亂的。}
\qaanswer{
普通話雖然語法混亂(這也能看出漢語使用者中其實是沒有什麼人文學菁英的，都近百年了，語法梳理還是如此混亂不堪，也沒有人在正經研究)，但因為使用人口的絕對巨大，很多使用方法甚至影響到了英語，大帝國的語言都有朝分析語發展的趨勢，所以她在全世界都是一門超強勢語言。但語言的強勢和語言的美(既精確又經濟)不能劃等號。
}

\chapter[可以請你簡單說明一下「成功神學」是什麼嗎？]{可以請你簡單說明一下「成功神學」是什麼嗎？}
\qaanswer{
這個您可以去Google，做伸手黨不如自力更生。
}

\chapter[我是之前那個問您老有無興趣和精力做影片的那位，其實我指的「比較系統的文法和發音影片」是「因為一般吳語母語者，沒有專業研究背景，最多隻熟悉自己的片區。但新吳越知識分子應該瞭解別的片區，而我本人確實也對此有很濃厚興趣。所以您老如果能出這種影片系列，最好不過，當然確實工程量較大，只是一個冒昧的詢問」。]{我是之前那個問您老有無興趣和精力做影片的那位，其實我指的「比較系統的文法和發音影片」是「因為一般吳語母語者，沒有專業研究背景，最多隻熟悉自己的片區。但新吳越知識分子應該瞭解別的片區，而我本人確實也對此有很濃厚興趣。所以您老如果能出這種影片系列，最好不過，當然確實工程量較大，只是一個冒昧的詢問」。}
\qaanswer{
影片承載的資訊量不如一篇條理清晰的論文。學習任何東西都需要支付成本。對一門還未形成規範和標準的語言做影片教學，工程量大，效果差，事倍功半。我能夠給出的論文定期會在我的Medium上更新。學不學在於吳語區人士自己的選擇。發音這種東西還沒有形成標準，吳語拼音這種音位拼音就是一個音位對應多個音值，這樣記錄時用同樣的拼寫方式吳語區各地的人按自己的習慣發自己地方的音即可。工具網站可以參考 http://wu-chinese.com/bbs/portal.php

另外據我現在所掌握的資訊，吳語區內即使是對吳語感興趣的吳語使用者多數也還是藍青吳語使用者。他們基本上都中了「漢字中心思維」和「漢語普通話思維」很深的「毒」，多數人習慣用「漢語普通話思維」來理解吳語建構，這個毒不排乾淨，學習以原生吳語建構的吳語播音體會造成學習障礙。學吳語播音體前最好把關於漢語語境下的語言學概念和學習習慣都去除掉。首先就從「正字」、「音韻」這兩點去除。吳語的漢字根據不同時代形成的詞彙有多種讀音，而多個讀音的詞往往用到的是同一個漢字，不需要弄出一些五花八門的生僻漢字。另外要區分實詞和虛詞，虛詞不要用漢字表述。
}

\chapter[可以說，一般人覺得中文的美，來自於不太理解很多詞彙的這種神秘主義嗎？]{可以說，一般人覺得中文的美，來自於不太理解很多詞彙的這種神秘主義嗎？}
\qaanswer{
我不是哲學家和詩人。什麼叫中文的美？如果是漢語文言文體系我覺得確實有古典美，因為他適合詩歌。但是對於現代漢語的美，抱歉，我體會不到。
}

\chapter[可以透漏您老 " 簡明吳漢辭典和文法 " 目前的工作進度和預估完成的進度嗎]{可以透漏您老 " 簡明吳漢辭典和文法 " 目前的工作進度和預估完成的進度嗎}
\qaanswer{
這個看窩老高不高興持續下去了。如果某一天窩老覺得吳語區的人普遍的德性實在不配擁有自己的語言，這個工作進行到一定時候所有資料會封存起來留給未來有德性的人承接下去。現在時間節點還未到，工作依然需要繼續。對於窩老本人也是考驗德性的一件事。一切看天命。
}

\chapter[吳語語法會比英語還困難嗎？]{吳語語法會比英語還困難嗎？}
\qaanswer{
吳語是一種話題性非常強的語言，一句話的焦點往往放在句子的最後。這門語言最關注的是事物狀態的變化，從這個方面來說某些情況下表述會比英語更為精確。但吳語不像英語這樣的「帝國語言」有大量來自殖民地的原語法體系外的不規則變化，而且像漢語族其他語言一樣，吳語沒有語態，沒有動詞變形，但和漢語族大多數語言一樣動詞有大量的體態和動結式但掌握起來不難。此外吳語有區分單複數，但只體現在領屬的變化上，很容易掌握。如果排除現代漢語為代表的官話語法的幹擾(主要指的是吳語區大城市人用的藍青話)沒有那麼多不規則的規範語法以外的用法。只有連讀變調這個語法對於非母語學習者來說會比較困惑。所以個人覺得吳語語法並不難於英語。
}

\chapter[現代漢語有多少詞彙和用法是直接借用吳語詞彙或間接受吳語影響？ 可以簡單舉一兩個例子嗎？]{現代漢語有多少詞彙和用法是直接借用吳語詞彙或間接受吳語影響？ \\[0.35em] 可以簡單舉一兩個例子嗎？}
\qaanswer{
我覺得是那些中古漢語中被官話語法體系接受，而且又屬於從中古漢語演化爾來的吳語詞彙或是工具性詞彙才被現代漢語普通話所使用。直接借用的比如有「尷尬」這類表達細膩情感的詞語和「會鈔」等金融術語。間接受吳語影響的詞彙就不太說的清楚了，這個各人有各人的看法。
}

\chapter[徽州再人口多點，有可能great again嗎?很期待上海帶領東吳，溫州帶領南吳，徽州帶領西吳的局面。]{徽州再人口多點，有可能great again嗎?很期待上海帶領東吳，溫州帶領南吳，徽州帶領西吳的局面。}
\qaanswer{
時代不同了。不過徽州一直有武德，另外吳越的各地是各自發展自己的特色，不需要誰來帶誰。
}

\chapter[如果語法並不能區別於官話的話，就算像香港、南粵一樣，很多字不用正確漢字用俗字、有粵語作為官方語言的香港，但只要香港的媒體和教育也推普，被殖民同化會非常快。而比如東南亞百越裔能熟練掌握多種語言，就是因為這些語言語法不同，代表的思維、場景不同。所以，如果吳語不能高階化，不能語法上跟官話隔絕，最終還是會被同化。 我的理解對嗎？]{如果語法並不能區別於官話的話，就算像香港、南粵一樣，很多字不用正確漢字用俗字、有粵語作為官方語言的香港，但只要香港的媒體和教育也推普，被殖民同化會非常快。而比如東南亞百越裔能熟練掌握多種語言，就是因為這些語言語法不同，代表的思維、場景不同。所以，如果吳語不能高階化，不能語法上跟官話隔絕，最終還是會被同化。 \\[0.35em] 我的理解對嗎？}
\qaanswer{
如果不能讓自己的母語言文一致，不能用書面語表達自己母語的語言思維。那麼就不會有標準出來。因為標準就是需要可以被量化、記錄、傳播的。任何一種沒有標準的規範化的語法和構詞法的口語是無法在強勢語言的滲透下存活的。單純透過情懷、社會呼籲和家庭堅守來保護吳語是蒼白和無力的。
}

\chapter[之前群裡發過的影片，裡面溫州瑞安女子說的形容完美的「Lou thei」大概是什麼漢字呢？]{之前群裡發過的影片，裡面溫州瑞安女子說的形容完美的「Lou thei」大概是什麼漢字呢？}
\qaanswer{
這我就不知道了，對溫州詞彙還不是很瞭解，相關的書看的還不多。
}

\chapter[請問目前簡明漢吳辭典的進度？未來會有吳吳辭典嗎？ 播音體主要的語法，預計多久完成呢？半年？三年？]{請問目前簡明漢吳辭典的進度？未來會有吳吳辭典嗎？ \\[0.35em] 播音體主要的語法，預計多久完成呢？半年？三年？}
\qaanswer{
簡明漢吳辭典需要設計一個檢索程式，這塊還沒有技術支援，需要程式設計師的協助。有意者可以直接DM我的推特@narihatoohfu。

播音體語法已經完成的差不多了，但是還沒有打算完全公開，現在需要募集學習者。
}

\chapter[吳語的人稱代詞有性別的區分嗎？你妳他她？]{吳語的人稱代詞有性別的區分嗎？你妳他她？}
\qaanswer{
吳語區的人稱代詞各地情況不同差異很大。總的來說，原生吳語不區分性別，但建構中的書面語吳語播音體中在第三人稱單數做為主語出現時，有考慮用不同漢字來表示性別。如用「zy佢」代表男性，「zy伊」代表女性。
}

\chapter[您老說過上海市區的藍青吳語程度很嚴重。那麼，請問上海市區等地的原生吳語建構應該怎麼做？]{您老說過上海市區的藍青吳語程度很嚴重。那麼，請問上海市區等地的原生吳語建構應該怎麼做？}
\qaanswer{
如果不想20年內徹底被漢語普通話吞了的話，那麼老老實實學播音體吧。藍青話是非常不穩定的，本身就是官話底子，在藍青話基礎上推普毫無阻力。
}

\chapter[成鳩じいさん，有沒有考慮過用吳語和吳越歷史人設、搭配二次元妹子，好好推廣吳語一波嗎？(草)]{成鳩じいさん，有沒有考慮過用吳語和吳越歷史人設、搭配二次元妹子，好好推廣吳語一波嗎？(草)}
\qaanswer{
妳們這些進步主義小將不要叫我じいさん。語法、詞彙、通用行文的格式等都沒構建好怎麼推廣吳語？妳們有幾個肯老老實實先把語法搞懂？說幾句藍青吳語趙家人可能還不管，但是妳們如果推真吳語，那是不是等著在牆內被趙彈？
}

\chapter[有想過在牆外備份吳語的資料嗎？怕被趙彈......]{有想過在牆外備份吳語的資料嗎？怕被趙彈......}
\qaanswer{
妳想表達什麼意思？
}

\chapter[閣下整理出吳語播音體欲以為吳語書面語，那閣下是如何看待那些並沒有那麼多時態等的吳語方言的呢？]{閣下整理出吳語播音體欲以為吳語書面語，那閣下是如何看待那些並沒有那麼多時態等的吳語方言的呢？}
\qaanswer{
幾乎所有的吳語亞方言都或多或少有絕對時態標記(時態)和相對時態標記(事態)。但吳語區的府城被官話和普通話滲透、漂變過深，很多方面已經藍青化和普化，自然就看不到時態的痕跡了。所以問題在於，吳語構建按官話，普通話的語法來合適還是回歸原生吳語本源合適？走前者的路線短期無需支付大的成本但最終會使吳語徹底消亡，走後者的路線成本大且艱辛，但能讓吳語作為一種獨立語言延續下去，最重要的是能將吳語語言思維這種獨特的思維方式延續下去。
}

\chapter[目前，很少見到你著墨於構詞法(造詞的原則)？ 閣下的文章幾乎都是以梳理語法為主。 雖然語法重要性很大， 但是構詞法、辨別原生吳語詞、保留非藍青的詞彙或句型， 重要程度不遜於語法呢?]{目前，很少見到你著墨於構詞法(造詞的原則)？ \\[0.35em] 閣下的文章幾乎都是以梳理語法為主。 \\[0.35em] 雖然語法重要性很大， \\[0.35em] 但是構詞法、辨別原生吳語詞、保留非藍青的詞彙或句型， \\[0.35em] 重要程度不遜於語法呢?}
\qaanswer{
詞彙需要編纂『漢吳辭典』，這件事即使有結果在當下的貴國也是無法公開傳播的。
}

\chapter[你在書面語那篇文中提到：「德語250年......外加網際網路的傳播可以讓新吳語在20年左右普及......」「吳語如果沒有新的發展就談不上維持和保護。未來，無法做到言文一致的語言將很快在AI時代消亡。」 閣下眼中，是不是在新時代與AI當中， 看到了復興吳語 建構書面吳語新的可能性？]{你在書面語那篇文中提到：「德語250年......外加網際網路的傳播可以讓新吳語在20年左右普及......」「吳語如果沒有新的發展就談不上維持和保護。未來，無法做到言文一致的語言將很快在AI時代消亡。」 \\[0.35em] 閣下眼中，是不是在新時代與AI當中， \\[0.35em] 看到了復興吳語 建構書面吳語新的可能性？}
\qaanswer{
吳語不發展，不進行語法整理和詞彙疏理、構建，不透過書面語定型，不出30年就會徹底滅絕。所謂的8000多萬吳語使用者只是一個30年前統計的虛數。現下除去已經逝去的高齡吳語使用者，和現在絕大多數中老年藍青話使用者外，真正還會原生吳語，掌握一定量吳語詞彙的人估計100萬都不會有了。AI時代只是在整理和構建吳語時能夠用到的更為高效、便捷的工具，語言構建這件事本身如果沒人做的話到2030年代這個時候，區別於漢語官話的原生吳語使用者可能就已經不存在了。
}

\chapter[請問小字與吳語拼音及某吳語方言音系的關系是 實際音系→吳語拼音→小字 還是 吳語拼音←實際音系→小字？ (盡管說吳拼是音位拼音，但由於不同地區音系差異，吳拼方案總會有區別，尤其是與內部一致性高的北吳所不同的南吳。)]{請問小字與吳語拼音及某吳語方言音系的關系是 \\[0.35em] 實際音系→吳語拼音→小字 \\[0.35em] 還是 \\[0.35em] 吳語拼音←實際音系→小字？ \\[0.35em] (盡管說吳拼是音位拼音，但由於不同地區音系差異，吳拼方案總會有區別，尤其是與內部一致性高的北吳所不同的南吳。)}
\qaanswer{
漢字以音位拼音為主，各地可以先按自己的發音習慣來，這樣7成以上就是一個體系的，音值不同沒關係。等有了一定書面語創作的量就看哪個地區的創作者多了。露出度高的地區的發音音值習慣就有競爭力。這個得看自然演化。如果書面語都出不來，吳語完球是無法避免的。
}

\chapter[閣下認為在國內(牆內)公開或半公開做吳語相關的哪些工作到什麼程度才會被趙彈打擊？]{閣下認為在國內(牆內)公開或半公開做吳語相關的哪些工作到什麼程度才會被趙彈打擊？}
\qaanswer{
妳試一試不就知道了？
}

\chapter[請問，先生Medium的文章更新還會繼續嗎？]{請問，先生Medium的文章更新還會繼續嗎？}
\qaanswer{
會繼續。主要看有沒有時間了。寫過的東西需要系統的校對整理後發上去。
}

\chapter[你們編撰的吳語辭典，只需要先蒐集1000-3000生活常用字，就足夠使用用了吧？但編著辭典困難的點不在於詞彙數量，而是為了尊重吳語底下各方言，避免全吳語區只接受某人群、某地域的用法，因此要儘可能蒐集、整理出「數個吳語底下方言」的詞彙、發音和用法。為嚴謹起見，還要考據詞彙的歷史變化，沒辦法只靠單純蒐集現在的詞彙、發音，就編成一部辭典。這意味著，乍看之下編撰常用3000個詞彙，實質上要理清30000個沒有書面紀錄、散落在田野的詞彙？]{你們編撰的吳語辭典，只需要先蒐集1000-3000生活常用字，就足夠使用用了吧？但編著辭典困難的點不在於詞彙數量，而是為了尊重吳語底下各方言，避免全吳語區只接受某人群、某地域的用法，因此要儘可能蒐集、整理出「數個吳語底下方言」的詞彙、發音和用法。為嚴謹起見，還要考據詞彙的歷史變化，沒辦法只靠單純蒐集現在的詞彙、發音，就編成一部辭典。這意味著，乍看之下編撰常用3000個詞彙，實質上要理清30000個沒有書面紀錄、散落在田野的詞彙？}
\qaanswer{
需要一再反復強調：吳語不是字中心思維，是詞中心思維。所以並不是著出常用吳語漢字這麼簡單的事。編輯最基礎的漢吳辭典是為了能讓吳語區的人在學習吳語通用語(吳語播音體語法)語法後能夠用吳語書寫文章。所以常用詞彙肯定不止3000個。從語法和諧的角度講，有些吳語原本沒有的現代概念的詞彙不能直接使用漢語普通話詞彙，因此還需要按吳語語法構詞或半構詞。既有詞彙主要是收集和甄選全吳越地區的包括徽語區詞彙。目前雖然有不少所謂的吳語辭典，但都是吳漢辭典，實際上對於使用者而言是無法當作工具書查詢吳語詞彙並用來書寫吳語文章的。所以編輯吳漢辭典的意義很大。
}

\chapter[請問如何提出申請，以便系統學習吳語播音體？]{請問如何提出申請，以便系統學習吳語播音體？}
\qaanswer{
註冊推特，與@narihatoohfu聯絡
}

\chapter[有想過多在另一個網路平臺放文章嗎？像是 方格子，那裡貌似免費。]{有想過多在另一個網路平臺放文章嗎？像是 方格子，那裡貌似免費。}
\qaanswer{
不打算放在牆內或任何漢語為基礎語言的平臺
}

\chapter[2020下半年之前會發布辭典嗎？]{2020下半年之前會發布辭典嗎？}
\qaanswer{
不會。這是個漫長的過程。
}

\chapter[將來吳語的漢字會沿用正體字、簡體字，還是日本的簡化漢字？吳小字混用漢字時，漢字筆劃不少。]{將來吳語的漢字會沿用正體字、簡體字，還是日本的簡化漢字？吳小字混用漢字時，漢字筆劃不少。}
\qaanswer{
還是先擔心吳語在未來的30年內會不會滅絕吧。看現在的上海，已經基本死透了。語法徹底Mandarin化。
}

\chapter[如果吳語註定滅亡，若干年後，在原吳語區的人 學習 復興吳語，不就和 學習 另一面新的外語 沒有差別嗎？ 換言之，復興吳語 也不一定只能被吳語區裔人士學習。 但復興吳語不是純人工的語言，畢竟你們梳理的語法和辭典，取自現實中既存的語法和詞彙(以及思維)。舊的原生吳語曾存在過，只是舊的吳語註定滅亡罷了。]{如果吳語註定滅亡，若干年後，在原吳語區的人 學習 復興吳語，不就和 學習 另一面新的外語 沒有差別嗎？ \\[0.35em] 換言之，復興吳語 也不一定只能被吳語區裔人士學習。 \\[0.35em] 但復興吳語不是純人工的語言，畢竟你們梳理的語法和辭典，取自現實中既存的語法和詞彙(以及思維)。舊的原生吳語曾存在過，只是舊的吳語註定滅亡罷了。}
\qaanswer{
舊的因為普通話的洗滌消亡了，所以以後還是用普通話吧。要增加點異域情調可以用吳語腔調說普通話。完美⋯
}

\chapter[先不論吳語未來如何， 請問你怎麼看 受到官話思維影響的諸語言， 未來會如何發展？官話本身未來又會如何？ 使用這些語言的人，在可預期的未來， 將因官話思維深受其害、而不自知？]{先不論吳語未來如何， \\[0.35em] 請問你怎麼看 受到官話思維影響的諸語言， \\[0.35em] 未來會如何發展？官話本身未來又會如何？ \\[0.35em] 使用這些語言的人，在可預期的未來， \\[0.35em] 將因官話思維深受其害、而不自知？}
\qaanswer{
首先官話和普通話不是同一個概念，儘管英語對於兩者的解釋都是Mandarin，它包含了官話與普通話。普通話是官話的2.0版，融入了更多英語語法的建構成分(儘管這種吸收和建構非常失敗)。Mandarin未來能發展到如何我無法預言，但我相信以這種語言為母語的人群在思維邏輯上以及思辨能力肯定是不如屈折語和黏著語母語的人群的。妳不可能讓一個對時態、語態、事態、格、屬、性毫無概念的人能在根本的思維層面上明白和運用這些，而這些又和邏輯思維、思辨力息息相關。簡單來說Mandarin到目前為止都沒有很好的整理出一套可以自圓其說的語法。所以以Mandarin為母語的文明中出現思維精神分裂的縫合怪非常的司空見慣。不借助其他邏輯思辨更強的語言作為思維工具，單憑Mandarin甚至連學習高等數學都會十分吃力，成才率也更低。從這個方面來說，大眾爭論不休的是否廢除漢字的討論其實倒是很次要的問題。東亞諸國書面語離開漢字是非常不現實的，但用全漢字作為書面語表述工具是會阻礙語法建構和思維發展的。
}

\chapter[漢字為什麼會導致理性思辨崩毀呢？當中的運作機理和前因後果是什麼？]{漢字為什麼會導致理性思辨崩毀呢？當中的運作機理和前因後果是什麼？}
\qaanswer{
漢字不會導致邏輯思辨出現障礙，但文字用純漢字表述會導致邏輯思辨的混亂。理由如下：

1. 漢字基本都有實意，當漢字作為虛詞成分出現時，它的基本實意會干擾解讀者理解它作為語法標記時的作用和意思。全漢字文字讓虛詞、實詞的區分和識別變得困難；

2.虛詞、實詞難以區別和辨識就會導致語法建構出現問題，難以形成自洽的語法體系，甚至讓很多漢語母語者以為漢語沒有語法；

3.漢字的延展性導致了大量的望文生義，特別是成語、習慣語的字面意思模糊，非常容易導致不同時代背景的解讀者朝著相反的方向理解，著名的例子有「七月流火」、「空穴來風」、「不孝有三、無後為大」。

以上這幾個問題相互交織疊加讓漢語的語言思維更是變得混亂不堪，直接導致漢語母語者多數會在邏輯思辨層面上出現障礙。
}

\chapter[像越南或韓國那樣採用全拼音文字，有哪些缺點？ 為什麼其他語言完全採用拼音是沒有問題的，(前)漢字圈諸語言卻不行？]{像越南或韓國那樣採用全拼音文字，有哪些缺點？ \\[0.35em] 為什麼其他語言完全採用拼音是沒有問題的，(前)漢字圈諸語言卻不行？}
\qaanswer{
漢字圈文化的語言都受到中古漢語、近古漢語強烈的影響與滲透。漢字詞彙大量地置換了這些地區的原生詞彙，漢字已經作為詞根的形式深度紮根與這些語言的詞彙中。而漢字詞根比較的簡練，導致重音律高，全拼音化的結果會造成同音異議詞變多，尤其是朝鮮語，拋開漢字後閱讀都會造成一定的障礙。再加上東亞語言詞彙的虛詞詞根部分特徵不明顯，全漢字導致虛實詞難以區分，全拼同樣會造成虛實詞難以區分。這一點採取實詞用漢字，虛詞用假名的日語構建是最好的。除此之外，漢字圈文化中詞性標記整理的相對好的是日語和朝鮮語，前者雖然有大量的訓讀但問題是音素太少，後者則是全文讀系統，這都導致全拼情況下記憶詞彙會變的困難。但這些其實都不是最嚴重的問題，真正的問題在於全拼後導致文化的割裂，今人無法很好的繼承祖先的文明傳承。大量文明資訊因不懂漢字而造成根源上的丟失。這一點朝鮮半島人和越南人深有體會，這兩地文字改革後一直沒有出現著名的文學作品也從側面反映了這一點。
}

\chapter[這樣理解可以嗎? 文讀就像是 源自 該時代官話 的外來語， 白讀就是 該時代既有的詞語。(白讀不等於固有詞) 而在東亞大陸上，各語言屢次遭受各朝代官話滲透， 前一個朝代的文讀，會成為後一個時代的白讀。 可想而知，一個漢字的讀音會受各朝代影響，層層累加。 這就是梳理構詞的困難點嗎？]{這樣理解可以嗎? \\[0.35em] 文讀就像是 源自 該時代官話 的外來語， \\[0.35em] 白讀就是 該時代既有的詞語。(白讀不等於固有詞) \\[0.35em] 而在東亞大陸上，各語言屢次遭受各朝代官話滲透， \\[0.35em] 前一個朝代的文讀，會成為後一個時代的白讀。 \\[0.35em] 可想而知，一個漢字的讀音會受各朝代影響，層層累加。 \\[0.35em] 這就是梳理構詞的困難點嗎？}
\qaanswer{
基本上就是這麼一回事
}

\chapter[推主你好。如果請你向人們推薦三位對你影響大的語言學家的話，你將會推薦哪些位呢？感謝]{推主你好。如果請你向人們推薦三位對你影響大的語言學家的話，你將會推薦哪些位呢？感謝}
\qaanswer{
只有Avram Noam Chomsky。東亞官方認可的知名的語言學家全都是渣，不知名的有一位，為了保護他，我不能透露他的姓名。
}

\chapter[你怎麼看待加拿大的心理學家 Jordan B Peterson ?]{你怎麼看待加拿大的心理學家 \\[0.35em] Jordan B Peterson ?}
\qaanswer{
不瞭解他，所以不敢黑屁
}

\chapter[吳語晚安怎麼說? 明朝見的意思無法和晚安對應。]{吳語晚安怎麼說? \\[0.35em] 明朝見的意思無法和晚安對應。}
\qaanswer{
各地沒有統一的說法。但各地原生吳語基本不用「時間+問候語」這樣的形式。府城受漢語官話和普通話影響有用「時間+問候語」的形式，但這不是吳語本有的表現。比如北部吳語區的府城有「日安」、「夜安」的說法，屬於明清官話體系的用法。現在也逐漸被「早上好」、「夜到好」這樣的偏普通話語法的句式用法所取代。

另外「明朝 見」是官話+普通話的混合用法，府城一般用「明朝 會」，「見」、「會」都是讀原調，與「明朝」並不是黏附關係，沒有變調，意思是「再見」。
}

\chapter[請教一下，如果上海話想一定程度擺脫藍青話的影響，用英語或日語作為替代是否可以短期內緩解危機，爭取一點時間呢？]{請教一下，如果上海話想一定程度擺脫藍青話的影響，用英語或日語作為替代是否可以短期內緩解危機，爭取一點時間呢？}
\qaanswer{
用英語和日語替代和官話/普通話給吳語帶來的藍青話本質上都會導致吳語成為皮欽話。因為他們的語法都與吳語語法體系不相容。解決危機的辦法從來都是在直面問題的態度下，對症下藥，而不是靠引入自己無法掌控的力量和因素來撞大運。基本上這種運氣永遠來不了。
}

\chapter[為什麼古代吳語被鮮卑人的語言間接影響不是壞事？ 語言變得更豐富， 而沒有變成像藍青話、洋徑濱或吳普這種皮欽語、克里奧語？]{為什麼古代吳語被鮮卑人的語言間接影響不是壞事？ \\[0.35em] 語言變得更豐富， \\[0.35em] 而沒有變成像藍青話、洋徑濱或吳普這種皮欽語、克里奧語？}
\qaanswer{
問題重新問一遍，字型統一用正體字吧。否則無法全部顯示。
}

\chapter[作為吳越本位主義者，您怎樣看待日中戦爭？]{作為吳越本位主義者，您怎樣看待日中戦爭？}
\qaanswer{
我是不是吳越本位主義者與我對近現代的中日戰爭的看法沒有什麼關係。

中日戰爭的本質就是雙方爭奪遠東代言人地位，雙方都想重新為東亞、北亞、印度支那、東南亞的秩序定下一個自己可以有絕對話語權的基調。戰爭期間掌握和控制中國武裝力量的是KMT以蔣介石為首的一幫江浙籍幹部，當然他們的身分認同是中華漢人而不是吳越人。
}

\chapter[為什麼說元代的漢兒言語並不能僅看作是誕生於金元時代的晚近語言，而是有著更久遠的上游歷史？可否將元代漢兒言語理解為中古漢語隋唐五代西北方音口語的直系後代？]{為什麼說元代的漢兒言語並不能僅看作是誕生於金元時代的晚近語言，而是有著更久遠的上游歷史？可否將元代漢兒言語理解為中古漢語隋唐五代西北方音口語的直系後代？}
\qaanswer{
首先漢兒言語是中古漢語發展到末期仍然保留大量中古漢語口語語法特徵的語言。而中古漢語口語是西晉滅亡到明初燕王掃北間整整800多年東亞大陸使用的口語。它是鮮卑語，鮮卑語分化出來的北方諸遊牧民族語言和上古漢語不斷融合疊加再融合的產物。中華史觀的人把北亞、東北亞民族按時代割裂成一個個獨立的毫不相干的民族，但事實上自匈奴起，北亞、東北亞的遊牧民族始終是同一波人以及他們彼此融合的子孫。鮮卑當中有很多部落來自早前的匈奴部落，要有一些內亞雜胡部落，他們也不是講統一的語言，與來自匈奴系鮮卑同源的柔然是蒙古的祖先，所以有理由相信他們和在東亞的中原地區建立北魏、北周政權的更早來自匈奴部落的拓拔氏、余文氏所使用的鮮卑語是同源一脈的。可以理解為早一批講匈奴/鮮卑混合語的鮮卑人在漢化中造成了語言的融合，慢慢的成為一種在東亞口頭通用的克里奧爾語。而之後蒙元又在此基礎上二次融入，讓東亞北部的口語更加偏近匈奴/鮮卑一系的語言。值得指出的是在明代建立朱元璋頒布『洪武正正韻』後明政府系統的禁止漢兒言語的使用，特別是朱棣奪權掃北後造成現在的官話區的口語面貌完全看不到漢兒言語留下的痕跡，倒是六南的原生語比如吳語、贛語甚至原生粵語，言從語法、詞彙更偏近中古漢語口語，許多方面能找到漢兒言語的影子。
}

\chapter[那 洪武正韻(明朝的官話)又是源自於那裡? 有一種謬論是：現在的普通話(北京官話)是受滿州人影響的方言。 但其實現代普通話是源自於明朝官話， 明朝官話與鮮卑人或吳越人沒有太大的關係？]{那洪武正韻(明朝的官話)又是源自於那裡? \\[0.35em] 有一種謬論是：現在的普通話(北京官話)是受滿州人影響的方言。 \\[0.35em] 但其實現代普通話是源自於明朝官話， \\[0.35em] 明朝官話與鮮卑人或吳越人沒有太大的關係？}
\qaanswer{
這得問朱重八為何要將淮西流民區的口語和宋元南京話的音韻體系當作正統了。事實上按『廣韻』、『集韻』的脈絡來看，『洪武正韻』也是不按規矩來的胡鬧。流民軍建立的王朝文化上先天的畸形。明朝官話口語的語法和中古漢語口語出現了斷層，音韻上演化自中古漢語音韻，詞彙還是保留了很多中古漢語詞彙，吸收了各地非官話區和行業行會的詞彙。吳語區因為在南宋城市化達到頂峰，社會多樣性高，市民生活豐富多彩，湧現了大量的戲曲、小說等市民文學作品，創造了大量的詞彙。它們中的很多被明清官話吸收。
}

\chapter[請問吳語的副詞、副動詞的系統閣下有過闡述嗎？比如uli、suli、syli、deuli等等這些]{請問吳語的副詞、副動詞的系統閣下有過闡述嗎？比如uli、suli、syli、deuli等等這些\begin{goetsu}
    
    \end{goetsu}}
\qaanswer{
這些都是副詞/副動詞詞根，一般有以下幾種

時間副詞詞根：

ka(屆)

katsy(屆子)

deu(頭)

非時間副詞詞根：

ka(介)

nenka(然介)

li(裡)

uli(乎裡)

deuli(頭裡)

syli(勢裡)

syuli(勢乎裡)

suli(勢乎的合音+裡)

tsaonsyli(場勢裡)

konsyli(功勢裡)

san(生)*通常也作形容詞詞根用

kuae(關)

副動詞詞根：

tsy(子)
}

\chapter[新冠肺炎疫情的衝擊，牆內海外無一倖免， 這對吳語復興有很嚴重的影響了吧？ 時間線更加緊迫了。 祝早日順利建構好、階段性完成部分辭典， 祝現有的語言記憶和文物儲存良好。]{新冠肺炎疫情的衝擊，牆內海外無一倖免， \\[0.35em] 這對吳語復興有很嚴重的影響了吧？ \\[0.35em] 時間線更加緊迫了。 \\[0.35em] 祝早日順利建構好、階段性完成部分辭典， \\[0.35em] 祝現有的語言記憶和文物儲存良好。}
\qaanswer{
吳語播音體被世人嚴肅的對待，並被推廣學習極大可能是在吳語區現有吳語(包括吳音藍青話)消亡以後的事了。只有擁有文字紀錄的語言才有可能做為標本留存下去，而在未來，吳語播音體可能是唯一可信的證明吳語曾經存在過，並區別於Mandarin的證據了。而未來在吳語業已消亡的很長時間之後，如果還能復活，很有可能靠的是現在構建中的基於梳理原生吳語後的吳語播音體的體系。我們所要做的就是為未來保留這種可能性。
}

\chapter[當代-吳越民族， 是否要面對 五代宋-江淮類似的命運？ (作為東亞大陸社會相對健全的地區， 即將面對經濟文化不可逆的重傷？) 江淮盛衰期： 南北朝(-傷) 隋唐(++++極盛) 五代(-衰) 宋(- -再衰) 明(- - -大洪水) 明以後到清末 (- 自然與社會生態被京杭運河嚴重破壞) 清末 (- -運河失效) 吳越盛衰期： 隋唐(+小盛)- 五代(+++極盛) 宋元(++小盛) 明初(- -小洪水) 晚明 清初(+小盛) 太平天國(- - 小洪水) 晚清民初(+小盛) 抗戰(- 衰)- CCP初(- -小洪水) 改革開放-2018年 (++經濟大盛，文化衰) 中國衰退，臺海戰爭前線 (- - -潛在經濟、文化大衰退的可能？)]{當代-吳越民族， \\[0.35em] 是否要面對 五代宋-江淮類似的命運？ \\[0.35em] (作為東亞大陸社會相對健全的地區， \\[0.35em] 即將面對經濟文化不可逆的重傷？) \\[0.35em] 江淮盛衰期： \\[0.35em] 南北朝(-傷) \\[0.35em] 隋唐(++++極盛) \\[0.35em] 五代(-衰) \\[0.35em] 宋(- -再衰) \\[0.35em] 明(- - -大洪水) \\[0.35em] 明以後到清末 \\[0.35em] (- 自然與社會生態被京杭運河嚴重破壞) \\[0.35em] 清末 \\[0.35em] (- -運河失效) \\[0.35em] 吳越盛衰期： \\[0.35em] 隋唐(+小盛)- \\[0.35em] 五代(+++極盛) \\[0.35em] 宋元(++小盛) \\[0.35em] 明初(- -小洪水) \\[0.35em] 晚明 清初(+小盛) \\[0.35em] 太平天國(- - 小洪水) \\[0.35em] 晚清民初(+小盛) \\[0.35em] 抗戰(- 衰)- \\[0.35em] CCP初(- -小洪水) \\[0.35em] 改革開放-2018年 \\[0.35em] (++經濟大盛，文化衰) \\[0.35em] 中國衰退，臺海戰爭前線 \\[0.35em] (- - -潛在經濟、文化大衰退的可能？)}
\qaanswer{
如果保護不了自己的勝利果實，那之前的繁盛榮昌都是虛假的。
}

\chapter[請具體說明何為 文法去官話化？ 官話思維是什麼意思？]{請具體說明何為 文法去官話化？ \\[0.35em] 官話思維是什麼意思？}
\qaanswer{
語灋(文灋)就是一種語言的思維方式和思維邏輯，明白這一點，就明白我說的話了。
}

\chapter[推主可以用吳語播音體翻譯一遍東亞語言對比材料中的名篇《北風與太陽》嗎？這篇小故事覆蓋了好多語言現象。 附原文：「 有一回，北風跟太陽在那裡爭，誰的本事大。爭來爭去分不出高低。這時候路上來了個過路人，身上穿著件厚大衣。他們倆就講好了，誰有本事先叫這個過路人把他的大衣脫下來，就算誰的本事大。北風就拼命地吹起來，不過他越是吹得厲害，那個過路人就把大衣包得越緊。後來北風沒辦法了，只好拉倒了。過了一會兒，太陽出來了。他熱烘烘地一曬，那個過路人馬上就把那件大衣脫下來了。這下北風只好承認，他們倆當中還是太陽的本事大。」]{推主可以用吳語播音體翻譯一遍東亞語言對比材料中的名篇《北風與太陽》嗎？這篇小故事覆蓋了好多語言現象。 \\[0.35em] 附原文：「 有一回，北風跟太陽在那裡爭，誰的本事大。爭來爭去分不出高低。這時候路上來了個過路人，身上穿著件厚大衣。他們倆就講好了，誰有本事先叫這個過路人把他的大衣脫下來，就算誰的本事大。北風就拼命地吹起來，不過他越是吹得厲害，那個過路人就把大衣包得越緊。後來北風沒辦法了，只好拉倒了。過了一會兒，太陽出來了。他熱烘烘地一曬，那個過路人馬上就把那件大衣脫下來了。這下北風只好承認，他們倆當中還是太陽的本事大。」}
\qaanswer{
可以試一試
}

\chapter[你怎麼看待月絕書標準吳語的構建？]{你怎麼看待月絕書標準吳語的構建？}
\qaanswer{
不做評價。個人認為，吳語書面語的構建，一切得基於原生吳語的語言思維進行建構。脫離這一前提的構建或是便乘Mandarin語灋而抄襲拼湊而來的吳音藍青話只是現代漢語普通話更為分析化，語言思維邏輯更為混亂的拙劣仿製品，在現代漢語普通話的衝擊下基本不可能會有存在感，是毫無意義和價值的行為。
}

\chapter[吳語播音體有理性、感性這個詞嗎？ 另外，近現代從西方直接或間接引進的詞彙(更準確地說，新概念)， 是不是要重新造吳語詞彙了？]{吳語播音體有理性、感性這個詞嗎？ \\[0.35em] 另外，近現代從西方直接或間接引進的詞彙(更準確地說，新概念)， \\[0.35em] 是不是要重新造吳語詞彙了？}
\qaanswer{
不是吳語播音體有理性和感性對應的詞彙，而是吳語本身就有對應到理性和感性的對標詞彙。只要符合吳語構詞灋的規則，既有的已被現代漢語拼音話收錄的外來新概念詞彙都可以被吳語播音體接納。新構詞語只有在吳語既有詞庫中遍尋不著的前提下才會考慮直接從外來語吸收或是按照吳語構詞法新構詞彙。
}

\chapter[其實，早就有臺語(閩南語)或其他各東亞大陸語言的使用者， 零星地整理他們母語的語法。 但是，我引用推友 並明 所說的： 『用「漢語語法」的體系來整理語法， 而不是原生語法整理語法。 其實是打著「整理」旗號做藍青化的工整吧。』 整理語法易，避開官話思維整理語法難， 這是一個很大的陷阱？]{其實，早就有臺語(閩南語)或其他各東亞大陸語言的使用者， \\[0.35em] 零星地整理他們母語的語法。 \\[0.35em] 但是，我引用推友 並明 所說的： \\[0.35em] 『用「漢語語法」的體系來整理語法， \\[0.35em] 而不是原生語法整理語法。 \\[0.35em] 其實是打著「整理」旗號做藍青化的工整吧。』 \\[0.35em] 整理語法易，避開官話思維整理語法難， \\[0.35em] 這是一個很大的陷阱？}
\qaanswer{
整理語法不易，現代漢語普通話語法和東亞大陸多數地區的原生語言，包括官話區內的各亞方言，本身就存在差異，只是現代漢語普通話與這些方言、語言之間有些僅僅是量的差異，有些則是質的不同。以吳語為例，她是一種與漢語族Mandarin語支 下的各語言，語言思維(語灋)差異非常大的獨立語言。她雖然收到漢語族Mandarin語支的強烈影響，她的一些方言，如上海話，已經呈現和Mandarin語灋越來越趨同的現象，但她最核心的語言思維還是普遍存在於吳語區多數地區的中高齡使用者腦海中的。
}

\chapter[你什麼時候能把完整的語法寫出來啊，我已經迫不及待了]{你什麼時候能把完整的語法寫出來啊，我已經迫不及待了}
\qaanswer{
先把既出的學好吧
}

\chapter[小字的 文en 能eon 什麼分別]{小字的 文en 能eon 什麼分別}
\qaanswer{
沒看懂這個問題。
}
